\chapter{Memory}\label{cha:memory}

\section{The 64\,KB view and the 256\,MB behind it}
The 45GS10's program counter is 16 bits: code executes inside a 64\,KB
window. The 256\,MB of physical memory is reached three ways:
\begin{itemize}[leftmargin=1.4em]
\item \textbf{Flat pointers}: \verb|LDA [ptr],Z| and the Q forms read and
  write any byte of the 256\,MB directly. Data never needs banking.
\item \textbf{Bank registers} (\verb|$D600|, one per 8\,KB block): write a
  28-bit base and that block of the window shows that memory. Bytes 0--2 set
  the base; byte 3 switches (bit 7 = off).
\item \textbf{DMA} (\verb|$D200|): copy, fill, swap, line, triangle ---
  instant, physical addresses.
\end{itemize}

\section{The CPU view, unmapped}
\begin{center}
\small
\begin{tabular}{@{}l l >{\raggedright\arraybackslash}p{0.48\linewidth}@{}}
\toprule
\texttt{\$0000-\$01FF} & system & zero page and stack\\
\texttt{\$0200-\$07FF} & system & the ROM's data and C stack, and two bytes worth knowing (below)\\
\texttt{\$0800-\$CFFF} & user   & 50\,KB for programs, ROM-free at launch\\
\texttt{\$A000-\$BFFF} & ROM    & \emph{the sideways window; RAM underneath}\\
\texttt{\$C000-\$CFFF} & ROM    & \emph{RAM underneath, bankable}\\
\texttt{\$D000-\$DFFF} & I/O    & always I/O, whatever is banked\\
\texttt{\$E000-\$FEFF} & ROM    & \emph{RAM underneath, bankable}\\
\texttt{\$FF00-\$FFFF} & stub   & always ROM: system-call stub and vectors\\
\bottomrule
\end{tabular}
\end{center}

\noindent Two bytes of system state are fixed, so that programs can
rely on them:

\begin{description}[leftmargin=1.4em,style=nextline]
\item[\texttt{\$022E} --- a script is running] non-zero while \verb|EXEC|
  (and so \pth{/STARTUP.BAT}) is feeding the shell. A program that would
  wait for a key --- \verb|TYPE|'s ``more'', say --- should not, because
  there is nobody there to press it.
\item[\texttt{\$03FF} --- the result] the shell's result code: 0 for
  success. A program sets it to say it failed, which is what an RX
  script reads as \verb|RC|.
\end{description}

\section{Programs}
A \texttt{.prg} begins with two addresses, where it loads and where it
starts, and the shell honours both. The C programs of \pth{/SYSTEM/BIN}
load at \texttt{\$6000}, Mad Pascal's at \texttt{\$0800}, the two
BASICs at \texttt{\$7000}, and \verb|MONITOR| at \texttt{\$E000}, in the
RAM under the ROM --- so that the memory a monitor is there to look at,
\texttt{\$0800} to \texttt{\$CFFF}, is left exactly as it was.

\section{System calls}
The page \texttt{\$FF00} is always the ROM, whatever is banked, and its
jump table is the whole of the interface a program needs:

\begin{center}
\small
\begin{tabular}{@{}l l >{\raggedright\arraybackslash}p{0.56\linewidth}@{}}
\toprule
\texttt{\$FF80} & CHROUT & print the character in A \\
\texttt{\$FF83} & CHRIN  & wait for a key; it comes back in A \\
\texttt{\$FF86} & GETIN  & a key in A, or 0 if none is waiting \\
\texttt{\$FF89} & LOAD   & name at (\texttt{\$F0}), to \texttt{\$F2--\$F5}; status in A, size in \texttt{\$F6--\$F9} \\
\texttt{\$FF8C} & SAVE   & name at (\texttt{\$F0}), from \texttt{\$F2--\$F5}, length \texttt{\$F6--\$F9} \\
\texttt{\$FF8F} & SHELL  & run the line A/X points at, as if it were typed \\
\texttt{\$FF92} & VIDEO  & put the ROM's video mode and palette back after a program drew \\
\texttt{\$FF95} & ARGS   & what followed the program's name: (\texttt{\$F0}), length in A \\
\bottomrule
\end{tabular}
\end{center}

\noindent Anything handed to a system call must lie below
\texttt{\$A000}: during the call the ROM is banked in over
\texttt{\$A000--\$FFFF}, and a line kept up there would read as the ROM's
bytes. \verb|SHELL| copies the line into the shell's own buffer before
running it, so the program's copy is left alone.

\section{RAM under the ROM}
\verb|rom_out()| (or three pokes to the bank registers) banks blocks 5--7
onto the RAM beneath the ROM: a program then owns
\texttt{\$0800-\$CFFF} + \texttt{\$E000-\$FEFF} = 57\,KB. Every system call
still works, because the calls go through the \texttt{\$FF00} stub page,
which banks the ROM in and out around them.

\section{Sideways ROM}
The operating system is bigger than its half of the 64\,KB --- so, like
the BBC Micro before it, the machine grew \emph{sideways}. The ROM file
is a 24\,KB base image plus appended 8\,KB banks paged into the
\texttt{\$A000-\$BFFF} window: bank~0 is the base image itself, banks~1
and~2 hold the colder commands, and bank~3 has the line editor, its
history and the directory commands. Chapter~\ref{cha:commands} says which
bank every command lives in. The shell pages the right bank in around
each call, and up to sixteen banks --- 128\,KB of operating system ---
fit in the scheme.

The ROM is full, and it is kept that way on purpose: a command whose
work can be a program on the disk becomes one --- \verb|TYPE| and the
monitor did in September~2026 --- because a program costs nothing to
load (the disk is the host's, and on the K4510's own Linux it is in RAM)
and can be replaced without touching the ROM.

The calls go one way, and this is the rule to hold on to if you write
code of your own for a bank: code in a bank may call resident code as
freely as it likes, but resident code must never call into a bank. At
that address a different bank holds something else entirely, and the
dispatch is the only thing that knows which bank is in. A routine that
everything uses belongs in the base image, whatever it costs there.

\section{RAM under the I/O page}
A \texttt{MAP} of block 6 hides \texttt{\$D000-\$DFFF} and exposes RAM:
with the ROM also banked away a program owns a contiguous field from
\texttt{\$0800} to \texttt{\$FEFF} --- 61.75\,KB of the 64. The trick
that makes it safe is that \texttt{MAP} is an \emph{instruction}, not a
register: the program that hid the I/O can always ask for it back, so
there is no way to lock yourself out. (The far-call gate is unreachable
while block 6 is mapped; unmap before far calls.)

\section{Programs bigger than the window}
The \texttt{K4SG} executable format loads segments to any physical address
and the far-call gate (\verb|$DF00|) calls between them: a plain \verb|JSR|
into slot $n$ banks descriptor $n$'s code in, and the callee's \verb|RTS|
banks it back out. \verb|SEGDEMO| shows two overlays sharing one
address.
